\scanpage{123}
de $i$. \deleted{Question : la structure de positivité est elle \unsure{associée|analogie} à $i|_{\mu_2}$ ?} Donc, grâce à cette circonstance
spéciale, il y a des choix de \deleted{\unsure{polarisati|positivité}} \unsure{structure|structures}
polarisantes privilégiés \unsure{de|dans|du} $G$ : les
$i|_{\mu_2}$ avec $i$ dans la classe ci-dessus.
\margin{\diagram{\begin{tikzcd}[ampersand replacement=\&]G_{m\,\mathbb C}\arrow[r,dash,"i"]\& G_{\mathbb C}\\G_{m\,\mathbb C}\arrow[r,"i'"]\& G_{\mathbb C}\end{tikzcd}}}
Quel est l'arbitraire ? On constate que l'ens.
\inserted{des} $i-i'$ \deleted{\unsure{correspond|conjugués} : $v_{\mathfrak p}(\lambda)\,v_{\mathfrak p'}(\lambda)$} définit un
sous-groupe \uline{d'indice fini} de $\operatorname{Hom}(G_{m\,\mathbb C},G^\circ_{\mathbb C})$,
mais cela ne donne aucune indication sur la
\unsure{composante|comportement} des $i_0-i'_0$, où l'indice $0$ désigne le
\unsure{restr.|contr.} à $\mu_2$. Question semble secondaire.

\margin{$\vert\!\vert\!\vert$}En tous cas, on a \deleted{\unsure{des|deux}} \inserted{un} paquet de
$\mathbb R$-foncteurs fibres privilégiés, définis en termes des
différents $i$.

\begin{center}\rule{2.2cm}{0.3pt}\end{center}

Autre structure spécifique sur le $G$ \unsure{provenant|pour}
d'un corps fini : d'après la façon de
décrire $G_{\mathbb C}=D_{\mathbb C}(M)$, $M\subset L^*$, comme $L$ est
stable par conjugaison complexe, qui commute
à $\operatorname{Gal}(L/\mathbb Q)$, on trouve dans $M$ donc
$G$ une involution canonique, notée
\[
u\longmapsto u'
\]
Dans \deleted{\unsure{D}}$\operatorname{Ker}j^*$ c'est la symétrie, d'autre part sur $\operatorname{Im}j$ c'est
l'identité, ce qui suffit pratiquement à la déterminer. Sur
$\operatorname{Hom}(G_{m,\mathbb C},G_{\mathbb C})$ $u\mapsto u'$ induit la conjugaison complexe.
% archive sheet number: 122
